\documentclass[11pt]{article}
\usepackage[margin=2.2cm]{geometry}
\usepackage{xcolor}
\usepackage{colortbl}
\usepackage{booktabs}
\usepackage{longtable}
\usepackage{array}
\usepackage[T1]{fontenc}
\usepackage{hyperref}
\usepackage{enumitem}
\usepackage{fancyhdr}
\usepackage{titlesec}

\definecolor{included}{RGB}{209,240,211}
\definecolor{excluded}{RGB}{250,214,214}
\definecolor{darkgray}{RGB}{60,60,60}
\definecolor{linkblue}{RGB}{0,90,160}

\hypersetup{colorlinks=true, linkcolor=darkgray, urlcolor=linkblue, citecolor=darkgray}
\titleformat{\section}{\normalfont\Large\bfseries\color{darkgray}}{\thesection}{0.6em}{}
\titleformat{\subsection}{\normalfont\large\bfseries}{\thesubsection}{0.6em}{}

\pagestyle{fancy}
\fancyhf{}
\fancyhead[L]{\small\textit{Nu(Pr, Ra) dataset --- sources and provenance}}
\fancyhead[R]{\small\thepage}
\renewcommand{\headrulewidth}{0.4pt}

\newcommand{\legend}{%
  \noindent
  \colorbox{included}{\rule{0pt}{1.6ex}\hspace{1.6ex}} \ \textbf{Included} in \texttt{nu\_dataset\_model.csv} \qquad
  \colorbox{excluded}{\rule{0pt}{1.6ex}\hspace{1.6ex}} \ \textbf{Excluded} (see Comment column)
  \vspace{0.5em}
}

\title{\textbf{Nu(Pr, Ra) Dataset: Sources and Provenance}\\[0.3em]
\large Machine-learning prediction of the Nusselt number in 2D Rayleigh-B\'enard convection}
\author{Compiled by Manya, for the project with Shreshthi and Prof.\ Ambrish Pandey}
\date{\today}

\begin{document}
\maketitle

\section*{Overview}
This document records exactly where every number in
\texttt{data/nu\_dataset\_model.csv} came from. Section~3.3 of Shreshthi \&
Pandey's paper (\emph{Turbulent Prandtl number in the near-wall region of
thermal convection}) trains a random forest regressor to predict the Nusselt
number $Nu$ from the Prandtl number $Pr$ and Rayleigh number $Ra$, using 55
data points drawn from five cited papers plus the paper's own simulations.
This document reproduces the relevant tables from the four sources we have
obtained so far, highlights which rows were kept, and comments on why the
rest were left out.

\legend

\section*{References}
\begin{enumerate}[leftmargin=1.6em, itemsep=0.6em]
\item Shreshthi and A.\ Pandey. \emph{Turbulent Prandtl number in the
near-wall region of thermal convection.} Project paper (\texttt{papers/shreshthi\_pandey\_2026\_turbulent\_prandtl\_number.pdf}).
\item A.\ Pandey. \emph{Thermal boundary layer structure in
low-Prandtl-number turbulent convection.} J.\ Fluid Mech., 910:A13, 2021.
\href{https://doi.org/10.1017/jfm.2020.961}{doi:10.1017/jfm.2020.961} ---
\href{https://arxiv.org/abs/2101.03051}{arXiv:2101.03051}.
\item A.\ Pandey and K.\,R.\ Sreenivasan. \emph{Transient and steady
convection in two dimensions.} J.\ Fluid Mech., 1015:A42, 2025.
\href{https://doi.org/10.1017/jfm.2025.10357}{doi:10.1017/jfm.2025.10357} ---
\href{https://arxiv.org/abs/2503.03080}{arXiv:2503.03080}.
\item A.\ Pandey, H.\ Tiwari, and K.\,R.\ Sreenivasan. \emph{Thermal
convection in 1, 2, 3, and 4 dimensions.} J.\ Fluid Mech., in press, 2026.
\href{https://arxiv.org/abs/2607.26372}{arXiv:2607.26372}.
\item E.\,P.\ van der Poel, R.\,J.\,A.\,M.\ Stevens, and D.\ Lohse.
\emph{Comparison between two- and three-dimensional Rayleigh-B\'enard
convection.} J.\ Fluid Mech., 736:177--194, 2013.
\href{https://doi.org/10.1017/jfm.2013.488}{doi:10.1017/jfm.2013.488}.
\textbf{\color{excluded!60!black}Not yet obtained} --- reports $Nu$ only as
scatter plots (Fig.~5a, 6b), not a table; paywalled, needs institutional
access to digitize.
\item Y.\ Zhang and Q.\ Zhou. \emph{Low-Prandtl-number effects on
global and local statistics in two-dimensional Rayleigh-B\'enard
convection.} Phys.\ Fluids, 36(1):015107, 2024.
\href{https://doi.org/10.1063/5.0175011}{doi:10.1063/5.0175011}.
\textbf{\color{excluded!60!black}Not yet obtained} --- fully paywalled, no
open copy found; needs institutional access.
\end{enumerate}

\noindent All four obtained sources describe the \emph{same} physical
setup: a closed 2D square box ($\Gamma=1$), no-slip walls, isothermal
top/bottom plates, adiabatic sidewalls, solved with the same spectral-element
code. This is confirmed empirically, not just from the papers' text: several
rows below report the exact same $Nu$ to the decimal at the same $Pr, Ra$ in
two different papers --- only possible if it is literally the same
simulation, reused across papers.

\section{Present paper --- Shreshthi \& Pandey, Table 1}
Closed 2D box, $Ra=10^9$ fixed, $Pr$ varied. All 13 rows are the paper's own
flagship simulations and are included.

\begin{longtable}{l l l l}
\caption{Present paper, Table 1 ($Ra=10^9$).}\label{tab:present}\\
\toprule
\textbf{Pr} & \textbf{Ra} & \textbf{Nu} & \textbf{Comment} \\
\midrule
\endfirsthead
\toprule
\textbf{Pr} & \textbf{Ra} & \textbf{Nu} & \textbf{Comment} \\
\midrule
\endhead
\bottomrule
\endfoot
\rowcolor{included} 0.021 & $10^{9}$ & 37.43 & Included \\
\rowcolor{included} 0.05 & $10^{9}$ & 41.51 & Included \\
\rowcolor{included} 0.1 & $10^{9}$ & 42.72 & Included \\
\rowcolor{included} 0.2 & $10^{9}$ & 45.98 & Included \\
\rowcolor{included} 0.35 & $10^{9}$ & 48.1 & Included \\
\rowcolor{included} 0.5 & $10^{9}$ & 48.64 & Included \\
\rowcolor{included} 0.7 & $10^{9}$ & 49.22 & Included \\
\rowcolor{included} 1 & $10^{9}$ & 50.87 & Included \\
\rowcolor{included} 2 & $10^{9}$ & 51.09 & Included \\
\rowcolor{included} 4.38 & $10^{9}$ & 52.67 & Included \\
\rowcolor{included} 7 & $10^{9}$ & 52.27 & Included \\
\rowcolor{included} 20 & $10^{9}$ & 53.41 & Included \\
\rowcolor{included} 50 & $10^{9}$ & 52.12 & Included \\
\end{longtable}
\section{Pandey (2021), Table 1}
Closed 2D box, $Pr=0.021$ fixed, $Ra$ varied. Runs labelled with an
`a' (e.g.\ 3a) are resolution checks of the preceding run --- the same
physical point, re-run at finer mesh, to confirm numerical convergence.
Only the base run is kept; its resolution-check twin is excluded since it
is not an independent data point.

\begin{longtable}{l l l l}
\caption{Pandey (2021), Table 1 ($Pr=0.021$).}\label{tab:pandey2021}\\
\toprule
\textbf{Run} & \textbf{Ra} & \textbf{Nu} & \textbf{Comment} \\
\midrule
\endfirsthead
\toprule
\textbf{Run} & \textbf{Ra} & \textbf{Nu} & \textbf{Comment} \\
\midrule
\endhead
\bottomrule
\endfoot
\rowcolor{included} 1 & $5\times10^{5}$ & 5.2 & Included \\
\rowcolor{excluded} 1a & $5\times10^{5}$ & 5.2 & Resolution check of run 1 \\
\rowcolor{included} 2 & $10^{6}$ & 6.3 & Included \\
\rowcolor{excluded} 2a & $10^{6}$ & 6.3 & Resolution check of run 2 \\
\rowcolor{included} 3 & $10^{7}$ & 10.8 & Included \\
\rowcolor{excluded} 3a & $10^{7}$ & 11.1 & Resolution check of run 3 \\
\rowcolor{included} 4 & $10^{8}$ & 19.6 & Included \\
\rowcolor{included} 5 & $10^{9}$ & 35.4 & Included \\
\rowcolor{excluded} 5a & $10^{9}$ & 36.2 & Resolution check of run 5 \\
\end{longtable}
\section{Pandey, Tiwari \& Sreenivasan (2026, in press)}
\subsection*{Table 1 --- closed 2D box, $Pr=0.021$ and $Pr=0.05$}
The $Ra=10^9$ rows for both $Pr$ values report exactly the same $Nu$
as the present paper's own Table 1 (37.43 and 41.51 respectively) --- the
same simulation, reused across papers. Excluded here to avoid counting it
twice.

\begin{longtable}{l l l l}
\caption{Pandey, Tiwari \& Sreenivasan (2026), Table 1.}\label{tab:pandey2026t1}\\
\toprule
\textbf{Pr} & \textbf{Ra} & \textbf{Nu} & \textbf{Comment} \\
\midrule
\endfirsthead
\toprule
\textbf{Pr} & \textbf{Ra} & \textbf{Nu} & \textbf{Comment} \\
\midrule
\endhead
\bottomrule
\endfoot
\rowcolor{included} 0.021 & $10^{7}$ & 10.97 & Included \\
\rowcolor{included} 0.021 & $2\times10^{7}$ & 13.04 & Included \\
\rowcolor{included} 0.021 & $3\times10^{7}$ & 14.42 & Included \\
\rowcolor{included} 0.021 & $5\times10^{7}$ & 16.58 & Included \\
\rowcolor{included} 0.021 & $10^{8}$ & 19.68 & Included \\
\rowcolor{included} 0.021 & $4\times10^{8}$ & 28.32 & Included \\
\rowcolor{excluded} 0.021 & $10^{9}$ & 37.43 & Matches Present paper exactly -- duplicate \\
\rowcolor{included} 0.05 & $10^{7}$ & 12.41 & Included \\
\rowcolor{included} 0.05 & $2\times10^{7}$ & 14.52 & Included \\
\rowcolor{included} 0.05 & $3\times10^{7}$ & 16.06 & Included \\
\rowcolor{included} 0.05 & $10^{8}$ & 21.64 & Included \\
\rowcolor{included} 0.05 & $3\times10^{8}$ & 29.27 & Included \\
\rowcolor{included} 0.05 & $5\times10^{8}$ & 33.5 & Included \\
\rowcolor{excluded} 0.05 & $10^{9}$ & 41.51 & Matches Present paper exactly -- duplicate \\
\end{longtable}
\subsection*{Table 2 --- ``RBC-P'', $Pr=1$, $D=2$ column}
This table uses a different solver (DHARA, finite-difference, versus
the spectral-element code used everywhere else) for a 1D/2D/3D/4D
dimension-comparison study. It is a \emph{different physical configuration}
from the closed box used in every other table here: at $Pr=1, Ra=10^8$ it
reports $Nu=13.6$, versus $Nu=25.01$ in Pandey \& Sreenivasan (2025)'s Table
1 at the identical $Pr, Ra$ (Table~\ref{tab:ps2025}). That gap is far too
large to be numerical noise. \textbf{All rows here are excluded} from the
training set.

\begin{longtable}{l l l}
\caption{Pandey, Tiwari \& Sreenivasan (2026), Table 2 (RBC-P, $D=2$, $Pr=1$).}\label{tab:pandey2026t2}\\
\toprule
\textbf{Ra} & \textbf{Nu} & \textbf{Comment} \\
\midrule
\endfirsthead
\toprule
\textbf{Ra} & \textbf{Nu} & \textbf{Comment} \\
\midrule
\endhead
\bottomrule
\endfoot
\rowcolor{excluded} $2\times10^{5}$ & 4.8 & Different configuration (RBC-P) -- excluded \\
\rowcolor{excluded} $5\times10^{5}$ & 6.2 & Different configuration (RBC-P) -- excluded \\
\rowcolor{excluded} $10^{6}$ & 6.9 & Different configuration (RBC-P) -- excluded \\
\rowcolor{excluded} $2\times10^{6}$ & 8.5 & Different configuration (RBC-P) -- excluded \\
\rowcolor{excluded} $5\times10^{6}$ & 10.3 & Different configuration (RBC-P) -- excluded \\
\rowcolor{excluded} $10^{7}$ & 11.8 & Different configuration (RBC-P) -- excluded \\
\rowcolor{excluded} $2\times10^{7}$ & 13.1 & Different configuration (RBC-P) -- excluded \\
\rowcolor{excluded} $5\times10^{7}$ & 13.3 & Different configuration (RBC-P) -- excluded \\
\rowcolor{excluded} $10^{8}$ & 13.6 & Different configuration (RBC-P) -- excluded \\
\rowcolor{excluded} $2\times10^{8}$ & 15.5 & Different configuration (RBC-P) -- excluded \\
\rowcolor{excluded} $5\times10^{8}$ & 20.4 & Different configuration (RBC-P) -- excluded \\
\rowcolor{excluded} $10^{9}$ & 24.6 & Different configuration (RBC-P) -- excluded \\
\end{longtable}
\section{Pandey \& Sreenivasan (2025), Table 1}
Closed 2D box, $Pr=0.1$ and $Pr=1$, wide $Ra$ range. As with the 2026
paper, the $Ra=10^9$ rows duplicate the present paper's own simulations
exactly (42.72 and 50.87) and are excluded.

\begin{longtable}{l l l l}
\caption{Pandey \& Sreenivasan (2025), Table 1.}\label{tab:ps2025}\\
\toprule
\textbf{Pr} & \textbf{Ra} & \textbf{Nu} & \textbf{Comment} \\
\midrule
\endfirsthead
\toprule
\textbf{Pr} & \textbf{Ra} & \textbf{Nu} & \textbf{Comment} \\
\midrule
\endhead
\bottomrule
\endfoot
\rowcolor{included} 0.1 & $10^{6}$ & 5.91 & Included \\
\rowcolor{included} 0.1 & $2\times10^{6}$ & 7.05 & Included \\
\rowcolor{included} 0.1 & $3\times10^{6}$ & 7.88 & Included \\
\rowcolor{included} 0.1 & $6\times10^{6}$ & 10.04 & Included \\
\rowcolor{included} 0.1 & $10^{7}$ & 11.77 & Included \\
\rowcolor{included} 0.1 & $2\times10^{7}$ & 15.35 & Included \\
\rowcolor{included} 0.1 & $3\times10^{7}$ & 17.5 & Included \\
\rowcolor{included} 0.1 & $10^{8}$ & 23.01 & Included \\
\rowcolor{included} 0.1 & $3\times10^{8}$ & 31.06 & Included \\
\rowcolor{excluded} 0.1 & $10^{9}$ & 42.72 & Matches Present paper exactly -- duplicate \\
\rowcolor{included} 0.1 & $3\times10^{9}$ & 59.15 & Included \\
\rowcolor{included} 0.1 & $10^{10}$ & 84.74 & Included \\
\rowcolor{included} 1 & $3\times10^{6}$ & 6.01 & Included \\
\rowcolor{included} 1 & $6\times10^{6}$ & 7.11 & Included \\
\rowcolor{included} 1 & $8\times10^{6}$ & 8.18 & Included \\
\rowcolor{included} 1 & $9\times10^{6}$ & 8.64 & Included \\
\rowcolor{included} 1 & $10^{7}$ & 9.03 & Included \\
\rowcolor{included} 1 & $1.5\times10^{7}$ & 13.33 & Included \\
\rowcolor{included} 1 & $2\times10^{7}$ & 14.75 & Included \\
\rowcolor{included} 1 & $3\times10^{7}$ & 17.27 & Included \\
\rowcolor{included} 1 & $6\times10^{7}$ & 21.48 & Included \\
\rowcolor{included} 1 & $10^{8}$ & 25.01 & Included \\
\rowcolor{included} 1 & $3\times10^{8}$ & 35.61 & Included \\
\rowcolor{included} 1 & $6\times10^{8}$ & 44.15 & Included \\
\rowcolor{excluded} 1 & $10^{9}$ & 50.87 & Matches Present paper exactly -- duplicate \\
\rowcolor{included} 1 & $3\times10^{9}$ & 67.51 & Included \\
\rowcolor{included} 1 & $10^{10}$ & 94.52 & Included \\
\rowcolor{included} 1 & $3\times10^{10}$ & 130.2 & Included \\
\rowcolor{included} 1 & $10^{11}$ & 183.3 & Included \\
\rowcolor{included} 1 & $3\times10^{11}$ & 262.3 & Included \\
\rowcolor{included} 1 & $10^{12}$ & 394.2 & Included \\
\end{longtable}
\section{Summary}
\begin{longtable}{l r r r}
\caption{Row counts by source.}\label{tab:summary}\\
\toprule
\textbf{Source} & \textbf{Rows in table} & \textbf{Included} & \textbf{Excluded} \\
\midrule
\endfirsthead
\toprule
\textbf{Source} & \textbf{Rows in table} & \textbf{Included} & \textbf{Excluded} \\
\midrule
\endhead
\bottomrule
\endfoot
Present paper, Table 1 & 13 & 13 & 0 \\
Pandey (2021), Table 1 & 9 & 5 & 4 \\
Pandey, Tiwari \& Sreenivasan (2026), Table 1 & 14 & 12 & 2 \\
Pandey, Tiwari \& Sreenivasan (2026), Table 2 (RBC-P) & 12 & 0 & 12 \\
Pandey \& Sreenivasan (2025), Table 1 & 31 & 29 & 2 \\
\midrule
\textbf{Total} & \textbf{79} & \textbf{59} & \textbf{20} \\
\end{longtable}
\noindent The 59 included rows across all four sources make up
\texttt{data/nu\_dataset\_model.csv}. Once the two missing papers
(references 5 and 6) are obtained, their tables should be added here
following the same process: reproduce the table, confirm the physical
setup matches (closed box, same boundary conditions) before including any
row, and update the summary above.

\end{document}